% Body of the pre-specification appendix (review point #20).
% No \section line: main.tex supplies the heading and the label.
% Every date is re-derivable from git; the commands are in docs/paper2/PRESPEC-LEDGER.md.

A study that reports negative results, ablations built to explain them, and theory derived
after the fact owes the reader a way to tell those apart. This section is that ledger. Every
date below is read from the repository's history rather than recalled, and the commands that
re-derive each one are listed in \texttt{docs/\allowbreak paper2/\allowbreak PRESPEC-LEDGER.\allowbreak md}. Five labels are
used: \textbf{pre-specified} (a dated design document predicted it before any run),
\textbf{confirmatory} (ran as specified), \textbf{diagnostic} (added to explain or repair an
observed result), \textbf{exploratory ablation} (added after seeing a result, to test a
hypothesis the result suggested), and \textbf{post-hoc} (an analysis or a theorem chosen
after the data existed).

\paragraph{What was pre-specified.} The design document is dated 2026-07-06 and its
predictions were committed before the first run. Fourteen items are pre-specified and
confirmatory: the threshold law's expected shape; the reach mechanism as an explicit
\emph{go/no-go} gate on the whole study; the wall position as the rarity knob; the
gate-miss exactness check against $(1-r)^N$; the pervasive-error control arm that makes the
axis separation a separation; the smooth-bump contrast, with ``danger stays low'' predicted;
the smooth-learner probe; the full-spec synthesis control; the mode-absent conditional,
including the ${\approx}60\%$ rate; the $\varepsilon$-sensitivity sweep, which the spec
listed under \emph{Risks} as required; the pendulum as the second instrument; the Qwen
cross-family check; and the pendulum synthesis arm. The tightest of these is worth naming,
because it is the strongest pre-specification claim the paper can make: the spec's runbook
predicted the mode-present branch as a \emph{contingency} --- ``either outcome is a finding;
if it repairs, that becomes its own section'' --- so the repair result is a pre-registered
branch and not a post-hoc pivot. The first confirmation landed $2$\,h $40$\,min after the
prediction was committed.

\paragraph{What was added after the result it responds to.} Fifteen items, and the five the
review asks about specifically are all in this group. The PatchField2D instrument was named
in the spec as one of two options for a second instrument, with no predictions attached, and
its campaign was built eleven days later in response to a review of the 1D work. The
\textbf{square ablation} was added two to three days after the 2D repair failure it explains,
and it is a genuine falsification test of a hypothesis that was itself post-hoc: the
curvature reading was born of the result the ablation then excluded. The
\textbf{guided-prompt ablation} was added two days after the same failure, and its two cells
are the most expensive in the paper ($320$ calls each, $640$ of the campaign's $1959$). The
\textbf{disjoint seed blocks} and the \textbf{$50{,}000$-rollout dependence re-measurement}
are both repairs of errors a review found --- pooled intervals over shared samples, and an
unstated independence assumption --- eighteen and eight days after the results they correct.
The \textbf{sharp-plateau reward variant} was added nineteen days after the ``below random at
every knob'' claim it rescues. The distrust-region mitigation appears in no version of the
design document at all; the paper presents it as a contribution rather than as a prediction,
which is the correct framing. Two further items changed the paper's own reading of a result
rather than confirming it: the per-episode fence census turned a ``boundary-mapping
transient'' into ``a growing fraction of outright failures'', and the behavioural artifact
audit shipped once with a vacuous check (a mis-keyed dictionary lookup meant its
partial-repair condition never fired), whose correction moved a headline count.

\paragraph{The theory is post-hoc, and part of it was explicitly out of scope.} Every
proposition-backing artifact in this paper was produced after all the data existed, most of
them because a measurement looked suspiciously exact and the question was whether something
could be proved. That is a legitimate way to find a theorem and an illegitimate way to claim
a prediction, so none of the theory is presented as predicted. One item is more than
post-hoc: the design document said of exactly this material, ``covering-number analogues are
open; mention as a limitation, do not attempt for paper 2''. The paper now contains four
such results. They were added late, and the first version of the fencing bound was wrong in
the direction that mattered --- a covering number where a packing number was needed, with a
hand-computed constant that was also wrong --- caught by a reader's question rather than by
the numeric audit, and corrected the same day.

\paragraph{The manuscript's own claims of pre-registration and falsification.} Ten phrases
in this paper assert something about our process. Seven are supported, one is supported as a
statistical rather than a temporal claim (the coverage certificate's cell family is fixed
before the data by \emph{construction}, which is what that argument needs and does not
require a dated record), and \textbf{two were not supported by any dated record} and have
been corrected in this revision:
\begin{itemize}
\item The sharp-plateau variant's competence risk was described as ``the pre-registered
risk''. No dated text records it before the result; the pre-run note for that variant states
a cost, a payoff and a different caveat. The text now says ``the risk recorded with the
variant's script''.
\item The 2D section said ``we predicted partial repair''. The instrumentation that detects
partial repair --- the per-mode sample flags and the four-way branch summary --- was written
in the same commit as the run, so the expectation is real and designed for; but the sentence
itself first appears one day \emph{after} the measurement. The text now says the instrument
was built to detect partial repair and that none occurred.
\end{itemize}
A third phrase, ``an axis-aligned square with flat edges excludes the curvature reading'',
is sound as logic but was silent about chronology; it now states that the ablation was added
after the result it responds to.

\paragraph{Summary.} Fourteen pre-specified and confirmatory items, all on the
one-dimensional story: the mechanism, the axis separation, the smooth contrast, both 1D
instruments, and the synthesis trichotomy including its contingency branch. Fifteen items
added after the results they respond to, of which five are diagnostic repairs driven by
review, five are exploratory ablations, three are post-hoc analyses and two are mixed. Ten
post-hoc theory results. The reader should treat the 1D synthesis result as confirmatory and
everything about the two-dimensional mode --- including the two ablations that exclude
curvature and prompting --- as exploratory work whose hypotheses were formed after the
finding.
